% New author continuation. No collision-normalized LLT is claimed here.
\section{The exact singular part of the physical roof derivative}
\label{sec:sbv-roof}

The complete roof variation in Sections~\ref{sec:av-all-labels} and
\ref{sec:wo-endpoint-tests} is a finite measure.  Finiteness alone does
not describe its singular part.  We now identify that part on the
\emph{entire} physical family, including every original section cut,
unvisited boundary and critical entrance.  The conclusion holds for
the actual smooth-guard and $C^1$ trajectory tests of
\eqref{eq:wo-measures}; those functions are not assumed subanalytic.
The geometric decomposition is made before those weights are inserted.

The finite physical graph and its parameterized semialgebraicity are
inherited from historical native164, not new assertions about the
orbit model.  Historical native40, native92 and native93 already extract and
reconstruct finite-count singular germs for their specified subanalytic
insertions.  On that older constructible scalar subclass, their
finite analytic preparation together with the retained BV bound
already excludes a Cantor part; no new priority is claimed for that
consequence.  Historical native153 already proves the signed
projected physical boundary current on its selected Morse charts,
including the combination of coincident decision interfaces.  That
signed-current principle is also inherited.  The new statement here
is its full physical derivative decomposition for the current
smooth observation class, with its exact physical coefficients
and the complete collision-uniform variation budget.  No definability
of a new observation is silently added to the hypotheses.

Fix $m\ge1$ and $R\in[9/20,47/100]$.  Write
\[
 f_{\ell}=f^{w;a,b,\Psi}_{\ell,m,R},\qquad
 D=\mathcal V_m^w(a,b,\Psi),\qquad \ell=(n,k),
\]
with precisely the choices and norms in
Theorem~\ref{thm:wo-bv}.  In particular $w=1$ or the original
$W^\varepsilon$, with $\varepsilon>0$ fixed, and there is exactly
one factor $c_*^{-1}$.  There are finitely many attainable labels at
this fixed collision count.  They are not pooled below.

\subsection{Pushforward along an actual boundary side}

\begin{lemma}[A common finite family of physical roof-side maps]
\label{lem:sbv-geometric-sides}
Fix the true collision count $m\ge1$. There is a finite integer $N_m$
and, for each original radius $R\in I$, a finite set
$\mathcal Z_{m,R}\subset[m\tau_-,m\tau_+]$, with
\[
                         \#\mathcal Z_{m,R}\le N_m,
\]
having the following property. In the signed integration-by-parts
formula for any original exact-record roof density in
Theorems~\ref{thm:rv-exact-bv} and \ref{thm:av-all-bv}, the regular
physical boundary terms can be written as a finite sum
\[
       (t_\gamma)_*(b_\gamma(r)\,dr).
\]
Here $r$ is a rational coordinate on a Gauss-direction arc,
$b_\gamma\in L^1$ is the signed or complex \emph{actual} one-copy
flux, and $t_\gamma$ is the original physical roof on that side.
Each $t_\gamma$ is $C^1$ and either constant, with value in
$\mathcal Z_{m,R}$, or has derivative nowhere zero on its open
parameter interval. The finitely many discarded parameter points
have zero flux. The positive-visited critical incoming contributions
are also supported on $\mathcal Z_{m,R}$.

The same set works simultaneously for every attainable exact label,
the marked phases, either $w=1$ or the original smooth guard
$w=W^\varepsilon$, and the bounded $C^1$ endpoint tests for which
Theorem~\ref{thm:wo-bv} proves finite variation. Neither a positive
incidence cutoff nor a replacement section is used. The set is
geometric: it is independent of $\varepsilon$, of the particular
smooth function $\Theta$, and of these test functions or phases.
No estimate uniform in $m$ is asserted for $N_m$.
\end{lemma}

\begin{proof}
We spell out the physical construction because local analyticity
alone would not give the asserted finite family uniformly in $R$.
The finite semialgebraic first-hit graph already proved in the
preserved native module 164 is the input here. We recall its
defining variables, rather than add a symbolic replacement model.

The uniform physical horizon gives a finite set of possible
next-center increments
\[
 \mathcal A=\{z\in\Lambda\setminus\{0\}:
                         |z|\le\tau_++2R_++1\}.
\]
For a word $C_0=0,\ldots,C_m$ with increments in $\mathcal A$,
use real contact points $Q_i=C_i+Rn_i$, unit normals $n_i$,
unit velocities $v_i$, and real lengths $\tau_i$.
The equations are
\[
 Q_{i+1}=Q_i+\tau_i v_i,\qquad
 v_i=v_{i-1}-2(v_{i-1}\cdot n_i)n_i,\qquad
 v_{i-1}\cdot n_i<0
\]
at every visited reflection, together with the original outgoing
sign at the initial collision and
$\tau_-\le\tau_i\le\tau_+$. To impose the first-hit rule, exclude
an interior intersection of the \emph{actual} flight with every
candidate disk preceding its designated endpoint. A disk that can
meet that segment lies within $\tau_++2R_+$ of its starting
center. Its squared distance to the segment is determined by the
two endpoint cases and by the interior orthogonal-foot case.
These are finite polynomial inequalities after the positive
projection denominator is retained. The current disk is treated
with its actual outgoing sign. No ray is restarted at a line foot.

The same finite enlargement gives every candidate for the original
least third-disk clearance.  A nearest-neighbor center other than
the terminal center has distance at most $1+R$ from the initial
contact, so its distance from the entire segment is at most $1+R$.
The least third-disk \emph{clearance} is therefore at most $1$.
A minimizing center is consequently within $\tau_++2R_++1$
of the starting center.
Square roots are encoded by their nonnegative root and squared
equation. Thus all physical incidences and original third-disk
clearances are semialgebraic on this graph.

Use finitely many rational charts for each normal circle, with
one-copy ownership and the finitely many seam sides retained until
cancellation. For a fixed chart rotation the charts have the form
\[
 n(x)=\operatorname{Rot}_{\alpha_*}
       \left(\frac{1-x^2}{1+x^2},\frac{2x}{1+x^2}\right),
 \qquad |x|<1.
\]
The moving center normal of each original rectangle is
semialgebraic in $R$, because its half-angle is the specified unique
root of the original polynomial in $q_R$. An angular side test is
$n\cdot n_*>\cos h$, with its actual fixed half-width $h$,
and its companion momentum test is linear. Fixed real numbers
such as $\cos h_0$ are allowed as coefficients.
Retain $\eta_0=\eta_m=1$,
$\sum_{i=0}^{m-1}\eta_i=n$, and $C_m=k$ exactly.
There are finitely many words and labels at this fixed $m$.
All these statements concern the deterministic physical graph and
the original roof $L=\sum_i\tau_i$.

On a regular endpoint sheet, write
$q=(\alpha_0,\alpha_m)$, $F=L$, and
$g=\nabla_qF=R(-p_0,p_m)$. The exact endpoint and optical
identities make $F$, $g$, and their derivatives semialgebraic in
rational normal charts: differentiation in $\alpha$ is
multiplication of the chart derivative by $(1+x^2)/2$.
Away from $g=0$, adjoin the positive variable $|g|$ and use
finitely many rational charts
\[
        e=\frac g{|g|}
           =\operatorname{Rot}_{\theta_*}
             \left(\frac{1-r^2}{1+r^2},
                         \frac{2r}{1+r^2}\right)
\]
for its true Gauss direction. The map $q\mapsto(F,e)$ has
the nonzero Jacobian $\chi$ of the retained optical coarea proof.
Consequently its regular fibers are discrete; since they are
semialgebraic, they are finite. The same statement holds for
the endpoint projection on a fixed regular word, by its nonzero
optical twist. Finite semialgebraic cell decomposition therefore
provides finitely many genuine inverse sheets, retaining their
physical multiplicities. It does not assign a nonphysical orbit
to a missing sheet.

Now take the complete one-sided graph of these sheets in variables
$(R,r,t)$, retaining the side identities before any observation
weight is inserted. Every closure and boundary in this construction
is taken \emph{fiberwise at the fixed radius}, before forming the
parameterized family. For example, membership of $(r,t)$ in the
closure of the $R$-fiber is expressed by
\[
 \forall a>0\ \exists(r',t')\quad
 (R,r',t')\text{ belongs to that sheet},\qquad
 (r-r')^2+(t-t')^2<a^2.
\]
Quantifier elimination makes this fiberwise closure semialgebraic.
The fiberwise relative boundary is obtained by the corresponding
Boolean operations. Thus limits from another radius never create
physical sides at an exceptional radius. At every fixed $R$ the
regular image is open in $(r,t)$, so its boundary has dimension at
most one. These boundaries include the original unvisited
constraints, caps used in the positive majorant, every actual
section side at every collision index, and their finite
intersections. They are semialgebraic, because closure, projection
and finite Boolean operations preserve semialgebraicity.
At positive visited incidences the critical set has
$g=0$ and positive endpoint Hessian. Its points are isolated on
each fixed-$R$ sheet. This is a zero-dimensional semialgebraic
fiber and hence finite. Over each such critical point the
Gauss-direction closure records its actual incoming sectors at
the single value $t=F(q)$. We keep those sectors, including
unvisited-tangency corners; they are not declared null because
the physical source point is null.

Apply $C^1$ cylindrical cell decomposition with $R$ ordered
before $r$. It is applied to the finite geometric side graphs,
their relative closures, and their projection-critical sets.
A side with nonconstant Gauss projection is then expressed on
finitely many cells as $t=t_\gamma(R,r)$.
For a fixed radius, further refine by the semialgebraic sets
where $\partial_rt_\gamma$ is zero, positive, or negative.
On an interval in the zero set the fundamental theorem of
calculus makes $t_\gamma$ constant. On each remaining interval
the derivative never vanishes. A side whose Gauss projection
is constant has zero $d\theta$-flux on its regular part and
does not create an omitted charge. Its intersections and the
finitely many point cells also have zero such flux.
The constant values and the critical roofs form
$\mathcal Z_{m,R}$.

This use of finite decomposition is precisely the finite
definable-family assertion in the cell and smooth-cell
decomposition theorems, for example
\cite[Theorems 2.9--2.10 and 6.6--6.7]{CosteOM}.
Their hypotheses apply to the semialgebraic graphs just
constructed. The finite cylindrical decomposition has finitely
many cells across the entire original radius interval, including
its exceptional parameter cells. Thus the number of its
fixed-radius interval pieces, and hence the number of their
constant values and critical roofs, has a finite common upper
bound $N_m$. These theorems provide no upper bound independent
of the formula size, which grows with $m$.

The inherited physical integration-by-parts formula on a roof
interval is the interior derivative plus its signed left and
right traces. For a regular side in the displayed chart,
$d\theta=2\,dr/(1+r^2)$ up to its chosen orientation.
Accordingly $b_\gamma(r)$ is the signed trace, including the
difference of the actual two incident record weights, times
this coordinate factor. Coincident genuine arcs have their
incident values combined before assigning the flux.
The physical trace and section-charge theorems give its $L^1$
bound; they do not require a bounded trace at each parameter
endpoint. Artificial angular charts, length cuts and
homogeneity cuts cancel between their two incident copies
before this representation. At a visited grazing limit the
optical trace is zero by Lemma~\ref{lem:rbv-endpoints},
including the simultaneous critical limit at its stated scope.
Thus it creates no additional side measure. The positive
visited critical fan is the separately retained incoming
measure described above.

Finally none of the smooth observation weights changes a
geometric sheet or introduces a jump side. The original
$W^\varepsilon$ is continuous on each physical sheet,
including where two branches of a least clearance exchange;
its interior derivative is the one already controlled in
the guarded BV proof. A bounded $C^1$ endpoint test similarly
adds only its proved interior derivative and multiplies
existing traces. Exact labels and discrete twists are constant
on the original section chambers. They therefore change
$b_\gamma$ and the critical sector coefficients but not the
finite side maps or their constant values. This proves the
claimed common set and its scope. An arbitrary discontinuous
Borel trajectory selector does not have this conclusion.
\end{proof}

\begin{lemma}[An integrable weight on a geometric roof side]
\label{lem:sbv-arc}
Let $J\subset\mathbb R$ be an open interval, let
$\gamma:J\to\mathbb R$ be $C^1$, and let $b\in L^1(J;\mathbb C)$.
If $\gamma'$ is nowhere zero, then
\[
 \gamma_*(b(r)\,dr)=h(t)\,dt,\qquad
 h(t)=\one_{\gamma(J)}(t)
       \frac{b(\gamma^{-1}(t))}{|\gamma'(\gamma^{-1}(t))|},
 \qquad \|h\|_1=\|b\|_1.
\]
If $\gamma\equiv\tau$, the same pushforward is
$(\int_J b(r)\,dr)\delta_\tau$.  Consequently a finite collection
of the geometric sides in Lemma~\ref{lem:sbv-geometric-sides},
with arbitrary integrable weights, has a pushforward which is the
sum of an $L^1$ density and a finite atomic measure.  A zero value
of $\gamma'$ at an isolated subdivision point creates no atom.
\end{lemma}
\begin{proof}
On each compact subinterval of $J$, a continuous nowhere-vanishing
derivative has constant sign and a positive lower bound in absolute
value.  The inverse-function theorem and the usual substitution
formula therefore give the first identity there, including the
displayed absolute denominator.  Exhaust $J$ by compact subintervals.
For $|b|$ this passage is monotone; for $b$ it follows by domination
with $|b|$.  This proves the full identity and norm equality without
a lower bound for $|\gamma'|$ near an endpoint of $J$.
For a constant map the pushforward identity is its definition.
The geometric lemma provides a finite subdivision into these two
cases and finitely many points.  Points have zero $b(r)\,dr$ mass.
Summing the absolutely integrable pieces proves the assertion.
In particular, a nonconstant side may produce an unbounded but
integrable density near a turning roof; it does not produce a
singular-continuous measure there.
\end{proof}

\subsection{The signed physical coefficient, before absolute values}

Use the one-copy physical integration-by-parts construction of
Theorem~\ref{thm:wo-bv}.  On each open regular roof--angle chamber
the actual coefficient is
\[
       w(t,\theta) A(t,\theta) V(t,\theta)v_\ell(t,\theta),
 \qquad V=a(x)b(T_R^m x)e^{i\Psi_m(x)},
\]
where $v_\ell$ is the original section/record indicator and
$A=\omega/(4\pi c_*\chi)$.  This construction differentiates the
complete physical density.  It does not differentiate an arbitrary
additional Borel trajectory selector.

For each regular geometric side $\gamma$ in a rational Gauss-angle
chart, let $r\in J_\gamma$ be its parameter, let
$t_\gamma(r)$ be its actual roof, and let
$b_{\ell,\gamma}(r)\,dr$ be its \emph{signed} coefficient in that
integration-by-parts formula.  Explicitly, on a chamber interface
with increasing-roof traces $C_\ell^-$ and $C_\ell^+$,
\begin{equation}\label{eq:sbv-side-coefficient}
 b_{\ell,\gamma}(r)
   =\bigl(C_\ell^+(r)-C_\ell^-(r)\bigr)
                           |d\theta/dr|.
\end{equation}
On the exterior side of a physical source the incident trace
$C_\ell$ is zero; the resulting boundary coefficient need not vanish.
At a regular section cut the continuous factor $wAV$ multiplies
the actual difference of the label indicators; an internal visit
is not charged as two visits of the same label.  At an entrance to
a physical interval the coefficient is positive for $w=V=1$;
at an exit it is negative.  For complex tests the signed formula
still applies.  Artificial chart seams are canceled before
\eqref{eq:sbv-side-coefficient} is formed.  Coincident real
constraints describe one interface with its two actual traces,
not extra copies of that interface.  A side with constant Gauss
angle has zero angular charge and is omitted.

Let $\mathcal K_{m,R}$ be the one-copy positive-visited critical
germs of the retained physical action.  Write $t_x$ for a germ's
critical roof and $G_x$ for its endpoint Hessian.  In the coordinates
$y=G_x^{1/2}(q-q_x)$ let $\beta_{\ell,y}(x)$ be the angular fraction
of its actual physical tangent sectors with the exact record $\ell$
and both original endpoint tests.  It is zero when the record is
absent.  At that germ the smooth weights have the one-sided limits
$w_x,V_x$ along its designated physical word.  The incoming
critical coefficient is
\begin{equation}\label{eq:sbv-critical-coefficient}
 c_{\ell,x}=w_xV_x\,\mathfrak J(x)\beta_{\ell,y}(x),
 \qquad \mathfrak J(x)=\frac{\omega(x)}{2c_*\sqrt{\det G_x}}.
\end{equation}
Formula~\eqref{eq:sbv-critical-coefficient} uses the retained
one-sided sector theorem.  It does not declare an unvisited
tangency to be a regular collision.  A critical limit with a visited
grazing contact has zero coarea trace by
Lemma~\ref{lem:rbv-endpoints}, even when several constraints meet.
The values $V_x$ are limits of the actual smooth observations on
that same word.  No time-reversal invariance of the section or of
those observations is required.

Let $\mathcal S_{m,R}=\mathcal Z_{m,R}$ be the finite set supplied
by Lemma~\ref{lem:sbv-geometric-sides}: it contains the critical roofs and the roof values of all
constant-roof regular sides.  It depends only on the physical
geometry and the fixed exact section at this collision count,
not on $a,b,\Psi$, $\Theta$, or $\varepsilon$.  For
$\tau\in\mathcal S_{m,R}$ set
\begin{equation}\label{eq:sbv-atomic-coefficient}
 a_\ell(\tau)
  =\sum_{\substack{x\in\mathcal K_{m,R}\\t_x=\tau}}c_{\ell,x}
    +\sum_{\substack{\gamma:\ t_\gamma\equiv\tau}}
                             \int_{J_\gamma}b_{\ell,\gamma}(r)\,dr.
\end{equation}
All terms at the \emph{same physical roof and same exact record}
are added with their signs before an absolute value is taken.
Coincident roofs in distinct records are not combined.  A zero
value of $a_\ell(\tau)$ is allowed.  In particular critical
coefficients alone need not be the full atomic part.

\begin{theorem}[No singular-continuous derivative for the full smooth class]
\label{thm:sbv-physical}
For every specified physical parameter and test there are
$d_\ell\in L^1(\mathbb R;\mathbb C)$ such that, exactly as
finite distributional measures,
\begin{equation}\label{eq:sbv-decomposition}
             Df_\ell=d_\ell(t)\,dt
                     +\sum_{\tau\in\mathcal S_{m,R}}
                                     a_\ell(\tau)\delta_\tau.
\end{equation}
There is no singular-continuous part.  The density $d_\ell$ is the
sum of the actual regular interior derivative and the
one-dimensional coarea densities of all nonconstant physical roof
sides.  The atomic coefficients are exactly
\eqref{eq:sbv-atomic-coefficient}.  Moreover
\begin{equation}\label{eq:sbv-budget}
 \sum_\ell\left(\|d_\ell\|_1+
                   \sum_{\tau\in\mathcal S_{m,R}}|a_\ell(\tau)|\right)
          =\sum_\ell\|Df_\ell\|_{\rm TV}\le D.
\end{equation}
In particular the bound retains all words, all original labels and
the original conditioning, with no label-count factor.  The bound
is uniform in $m,R$ when the norms in Theorem~\ref{thm:wo-bv}
are uniformly bounded.  Neither the number of sides and atoms nor
their separation is claimed uniform as $m\to\infty$.

The canonical BV traces satisfy
\begin{equation}\label{eq:sbv-jump}
            f_\ell(\tau+)-f_\ell(\tau-)=a_\ell(\tau).
\end{equation}
On every interval avoiding $\mathcal S_{m,R}$ the density has an
absolutely continuous representative.  These statements do not
make the entire density differentiable at a jump, nor turn a
finite-count geometric assertion into a collision-normalized
pointwise estimate.
\end{theorem}
\begin{proof}
The proof of Theorem~\ref{thm:wo-bv} constructs $Df_\ell$ by
integrating the regular chamber derivative and all actual one-sided
charges.  Its full positive source bound is finite, including the
inverse-gradient multiplier and every endpoint/section charge.
Consequently it is legitimate to use that signed source identity,
and then project it to the physical roof, before taking variation.

The regular part is absolutely continuous in $t$: at fixed Gauss
angle its signed density is the actual $t$ derivative of $wAV$ on
each chamber.  Its full integral of the absolute value is finite
by that same proof.  Fubini therefore gives an $L^1$ roof density
for its sum.  This argument uses the actual regular derivative of
the smooth observations and guard and needs no definability of
those functions.  In particular the minimum in the least-clearance
definition is handled by its retained absolutely continuous
one-dimensional representative; an artificial nearest-center
switch is not a new physical endpoint.

The complete noncritical geometric side collection has the finite
description of Lemma~\ref{lem:sbv-geometric-sides}.  Its signed
weights satisfy
$\sum_{\ell,\gamma}\int|b_{\ell,\gamma}|<\infty$ by the complete
boundary accounting.  On a nonconstant monotone side apply
Lemma~\ref{lem:sbv-arc}; its projected derivative is absolutely
continuous, including the one-sided approach to a turning endpoint.
On a constant-roof side the same lemma gives precisely the second
sum in \eqref{eq:sbv-atomic-coefficient}.  Endpoints of those
regular sides have zero angular integration measure.  Their
critical limits have instead already been assigned to the
incoming critical construction; this avoids counting a critical
atom twice.

At a positive-visited critical germ the inherited local sector
coarea theorem and continuity of the weights give exactly
\eqref{eq:sbv-critical-coefficient}.  The positive endpoint Hessian makes the critical configurations
isolated on each regular sheet; the finite semialgebraic
description in Lemma~\ref{lem:sbv-geometric-sides} makes their
fixed-count family finite.
All visited-grazing traces are zero as actual one-sided traces;
collision-area nullness is not substituted for this fact.  The
exhaustive classification in Lemma~\ref{lem:rbv-endpoints} and
the exact section accounting in Theorem~\ref{thm:wo-bv} leave
no further source measure.  In particular no singular-continuous
measure can remain after these pushforwards.  This proves
\eqref{eq:sbv-decomposition} and its coefficient formula.

The absolutely continuous and atomic parts in
\eqref{eq:sbv-decomposition} are mutually singular.  The total
variation of their sum is therefore the sum of their variation
masses, also for complex measures.  After all same-roof
cancellations in \eqref{eq:sbv-atomic-coefficient}, this gives
the equality in \eqref{eq:sbv-budget}.  Its inequality is the
existing complete variation theorem, not a separate sum of
per-label constants.

Finally, for fixed $\ell$ define a representative by integrating
the right side of \eqref{eq:sbv-decomposition} from any point to
the left of its compact physical support.  It has that same
distributional derivative and equals $f_\ell$ almost everywhere:
their difference has zero distributional derivative and vanishes
on a left exterior interval.  The $L^1$ primitive is continuous,
and an atomic primitive has the jump equal to its coefficient.
This proves \eqref{eq:sbv-jump} and absolute continuity away
from the finite atom set.
\end{proof}

\subsection{The normalized jump constraint for the original raw theorem}

Take $w=1$, $V=1$, and the original density $p_{n,R}(k,m,t)$.
Let
\[
 E_m(M)=\sup_{R,n,k}\esssup_{t\in\mathcal C_{m,R,n,k}(M)}
       |m^2p_{n,R}(k,m,t)-\mathcal L_{m,R}(k_1,k_2,t,n)|.
\]
The full retained arithmetic reference is continuous in the physical
roof at fixed $m,R,n,k$; every continued near-unit branch remains
in it.  Let
\[
 \mathcal S_{m,R,\ell}^{\rm cen}(M)
    =\mathcal S_{m,R}\cap
                \operatorname{int}\mathcal C_{m,R,n,k}(M).
\]
Only this interior condition permits comparison with both traces
inside the displayed central domain; an isolated endpoint of that
domain is not silently treated as an interior roof.

\begin{proposition}[Necessary full-label central jump payment]
\label{prop:sbv-central-jumps}
For every finite collision count,
\begin{align}
 m^2|a_\ell(\tau)|&\le2E_m(M)
 \quad\bigl(\tau\in\mathcal S_{m,R,\ell}^{\rm cen}(M)\bigr),
                                                   \label{eq:sbv-jump-max}\\
 \sup_R m^2\sum_\ell\sum_{\tau\in
                      \mathcal S_{m,R,\ell}^{\rm cen}(M)}
                      |a_\ell(\tau)|^2
       &\le2D_{\rm all}^{1}E_m(M).                 \label{eq:sbv-jump-energy}
\end{align}
Hence the original raw theorem implies that both left sides tend
to zero as $m\to\infty$.  These conclusions are necessary tests
of that theorem.  Neither one is asserted sufficient for it.
\end{proposition}
\begin{proof}
Fix a radius, label and interior central roof $\tau$.  The definition
of the essential supremum bounds the normalized error at almost every
roof on each side of $\tau$ in a small central interval.  Choose
Lebesgue points approaching from those two sides.  The BV traces
exist, and the retained reference has the same continuous limit
from each side.  Passing to the two limits gives
\[
 |m^2p(\tau+)-\mathcal L(\tau)|\le E_m(M),\qquad
 |m^2p(\tau-)-\mathcal L(\tau)|\le E_m(M).
\]
Subtract and use \eqref{eq:sbv-jump} to prove
\eqref{eq:sbv-jump-max}.  This operation is performed at each
fixed radius and uses no common exceptional roof set across all
radii.

Multiply the pointwise coefficient bound by $|a_\ell(\tau)|$,
then sum over all the specified physical central atoms and labels.
By \eqref{eq:sbv-budget} with $w=V=1$ this gives
\[
 m^2\sum_{\ell,\tau:\,\mathrm{central}}|a_\ell(\tau)|^2
 \le2E_m(M)\sum_{\ell,\tau}|a_\ell(\tau)|
 \le2D_{\rm all}^{1}E_m(M).
\]
The same constant holds at every radius, proving the asserted
supremum.  There is no multiplication by the number of labels
or physical words.  If $E_m(M)$ tends to zero, the two conclusions
follow in the original collision order.
\end{proof}

\paragraph{Scope.}
The theorem gives a common signed physical formula on the full
current smooth observation class, whose weights need not be tame,
and ties all its singular derivative mass to actual constant-roof
sides and critical entrances.  It does not suppress its positive critical trace, negative exit
charge or sharp integrable turning density.  In particular the
complete normalized local derivative charge in
\eqref{eq:ui-local-charge-suffices} is still not proved small.
The original same-roof shortage in \eqref{eq:lc-final-shortage}
also retains its exact source, comparison pool and ordered limits.
All previous directed-current, paired-flux and critical-orbit
theorems remain available with their original stronger or
different conclusions.  The present structural result is not
presented as a rediscovery of their mechanisms.
